% 集中定义 EEG 信号波形模块
\providecommand{\eegDataCard}[3]{ % 节点名/标签/样式
  \node[#3] (#1) at (0, 0) {}; % 数据外框
  \node[font=\scriptsize] at ([yshift=2.4mm]#1.center) {#2}; % 数据标签
  \begin{scope}[shift={([yshift=-2mm]#1.center)}, scale=.62] % EEG波形位置
    \pic {eeg signal}; % EEG波形
  \end{scope}
}
%
\tikzset{
  % EEG白色背景样式
  eeg signal bg/.style={
    draw=gray!25, % 背景边框色
    fill=white, % 背景填充色
    rounded corners=1pt, % 背景轻微圆角
    line width=.25pt % 背景边框宽度
  },
  % EEG波形线样式
  eeg signal line/.style={
    draw=black!65,
    line width=.58pt,
    line cap=round,
    line join=round
  },
  % EEG信号波形
  pics/eeg signal/.style={
    code={ % 定义可复用EEG波形
      \path[eeg signal bg] (-.70, -.26) rectangle (.70, .30); % EEG白色背景
      \draw[eeg signal line]
        (-.62, .00) --
        (-.56, .01) --
        (-.50, -.02) --
        (-.45, .04) --
        (-.40, -.07) --
        (-.35, .05) --
        (-.30, -.02) --
        (-.25, -.15) --
        (-.21, .04) --
        (-.17, .18) --
        (-.13, .06) --
        (-.09, -.06) --
        (-.05, .02) --
        (.00, -.11) --
        (.05, -.04) --
        (.10, -.07) --
        (.16, .03) --
        (.21, -.03) --
        (.26, .08) --
        (.31, -.09) --
        (.36, .02) --
        (.42, -.04) --
        (.48, .03) --
        (.54, -.01) --
        (.58, .00) --
        (.62, .00); % EEG折线
    }
  }
}
